\chapter{Conclusions}
\label{chapter_conclusion}

\section{Current-Stage Conclusion}

This initial report establishes a focused foundation for an AI-powered business
analytics system whose first implementation and empirical validation target
Bangladeshi retail pharmacies. The framework may later be extended to other SME
verticals, but each extension requires its own data, constraints and validation.
The central problem is not merely
the absence of prediction models. It is the missing, trustworthy connection
between operational records, temporally valid analytics and a feasible owner
decision. A forecast has limited practical value if its source cannot be traced,
its test contains future information, or the result ignores cash, stock and lead
time. A customer-risk score is similarly incomplete when its probability is
uncalibrated or the proposed contact does not respect consent.

The evidence leads to four conclusions. Bangladesh-specific studies make
workflow fit, perceived value, cost and local technology use necessary design
considerations \cite{hoque2016ict,shahadat2023digital,hazra2021mfs}. Retail and
inventory research requires explicit treatment of intermittency, hierarchy,
forecast error and stock consequences
\cite{fildes2022retail,syntetos2005intermittent,hyndman2011hierarchical,
syntetos2010stock}. Leakage control, imbalance-aware evaluation, probability
calibration and understandable evidence are scientific safeguards
\cite{kaufman2012leakage,saito2015pr,niculescu2005calibration,lundberg2017shap}.
The resulting research gap is integrative: these principles have not yet been
evaluated together in a bounded Bangladesh-oriented SME workflow.

The proposed \system{} framework responds through a traceable sequence of
multi-pharmacy operational evidence, historical snapshots, comparison of a
transparent baseline, pooled pharmacy model and eligible shop-adapted model,
uncertainty and explanation, shop-specific constraint-aware action ranking,
owner response and later outcome monitoring. Constraints include purchasing
budget, stock, pack or minimum-order rules, supplier lead time, expiry, storage
conditions and consent. Its proposed contribution is this
governed sequence and its evaluation, not the invention of a new classifier or
an unsupported claim that a complex model will always perform best.

Continuous operation means that current data and recommendations can refresh as
the pharmacy uses the software, while drift and outcomes are monitored. It does
not mean uncontrolled retraining after every transaction. A replacement model
must be versioned and must pass chronological comparison with the deployed model
and baseline before promotion. The system remains business decision support: it
does not diagnose, prescribe or recommend medicine dosage.

The report therefore supplies a testable path into the next phase: ethics and
multi-pharmacy partner-data readiness, pharmacy schema and constraint inputs,
model freeze, held-out-pharmacy and untouched-period evaluation, usability tasks and
decision-impact analysis. It does not yet establish system effectiveness. Final
claims must follow the observed evidence, including any result in which a simple
baseline wins, explanations show no measurable advantage, or the available data
cannot support a research question.
